
Globally, groups are building increasingly powerful AI systems. These systems are becoming embedded into many aspects of our lives and society and, notably, are operating with more autonomy than previous technologies. \term{Large Language Models} (\term{LLMs}) are a class of AI systems that have shown remarkable capabilities in understanding and generating natural language, audio and visual content. With these advanced capabilities, LLMs are being used to act both on our behalf and alongside us. The application of these systems motivates alignment between the systems and how we would want them to behave. Given the early stages of AI development and its impact, effective action and direction now could have outsized downstream effects.

In this report I focus on the \term{alignment} of LLMs; that is the process todo....  I argue that current alignment methods lack democratically grounded processes and instead have perverse incentives that make alignment with public values/good unrealistic. I provide a project proposal to research and implement a method to align LLMs to public values and evaluate the success of this alignment in a way that is democratically grounded, transparent and fair.

\section{The Alignment Process of LLMs}

\term{AI alignment} is the field of study and practice that focuses on ensuring that AI systems' goals and behaviors are in line with what designers and end users intend. The field itself is relatively new, but draws on older ideas from control theory, ethics, policy, and governance.

LLMs are typically trained on vast amounts of text data using a \term{self-supervised learning} approach where they learn to predict the next word in a sentence. LLM alignment generally refers to the process of making their outputs not just coherent and contextually relevant, but also consistent with what users (and society) would want. For example, a counselling assistant that prioritizes immediate user approval could validate harmful rationalizations (e.g., minimizing alcohol dependence) rather than supporting long-term wellbeing.

Various methods have been proposed and implemented to help align LLMs to better match user intentions and values. These methods generally involve generating a response to a prompt and then using a feedback mechanism to adjust the model's parameters. The response dataset creates an \term{alignment target}, that in practice is whatever best serves the deploying organization's objectives (e.g., growth, retention, revenue). This is often at odds with end users' intentions and values.

\section{Problems with Alignment processes}
todo
Define the problem 

Have two problems
Value selection problem (we need them to be for the public)
Transparency problem (we need to know how it did the alignemtn and how well aligned it is).
Technical alignment problem (how you change the models to be aligned) (mention technical alignment problem)

\section{A way to move forward: Democratically Grounded Alignment}

To solve the alignment problem we need to not just improve methods of making LLMs aligned to a target, but also need ways to defensibly define a target and evaluation process. Current methods in the entire AI stack lack transparency, public participation and public accountability. This makes them in-defensible in a democratic context. Therefore, I argue to move forward we need to use more democratically grounded methods to define alignment targets and evaluate alignment success.

We can improve the definition of alignment targets by using public participation methods to gather a diverse set of values and preferences. This can then be used to create an alignment target that fairly represents the public's values. Current evaluation methods for alignment success are \scare{AI driven}, opaque and mainly measure short term response quality. Even if grounded in human feedback these methods don't judge models on real world situations or over longer time horizons.

These two ideas of fair alignment targets and robust evaluation provide a pathway towards having AI that can be built by the people, for the people. I argue for an approach that removes perverse profit incentives from the development and deployment of such transformative technology. Instead to be replaced with strong public institutions that develop and deploy AI in a way that is by definition aligned to public values and good.


\section{Report Overview}

This report is structured as a chapter and two appendices. The chapter is designed to provide the context needed to motivate and support the project proposal in the second appendix. The first appendix provides additional reading to either help inform the reader and expand on the main ideas.

\subsection*{Chapter 2: Background Study}

In this chapter I provide a comprehensive background study covering alignment as a field, technical LLM alignment, human values and a discussion of the sociopolitical context of AI. This provides the context for what is currently known in the field, what is wrong and how we might move towards a better future.

\subsection*{Appendix A: Related Ideas}

Here I introduce and discuss related ideas that are relevant to the main proposal yet not central to it. These include an overview of current and future AI safety concerns as well as an introduction to LLMs to help inform readers who are less familiar with the technical details.

\subsection*{Appendix B: Proposal}

Here I will propose a project to research and implement a method to align LLMs to public values and provide a strong democratically grounded evaluation methods to measure the success of alignment. This fills the gap between current alignment methods and the need for scalable, democratically grounded alignment techniques.